% Autorska priprava k Realisticky.cz 020801–020802; ulohy jsou autorske a obmenene.
% Zamerne nenahrazuje vyklad presnymi kopiemi cizich textu/prikladu.
\documentclass[11pt,a4paper]{article}
\input{../../sablony/pripravy/priprava.sty}
\usepackage{tikz}
\nastavPripravu{1. ročník}{Rovnice a nerovnice}{90}{Lineární rovnice s parametrem}
\begin{document}
\blok{Návaznost a motivace}
Od běžné rovnice se parametrická liší tím, že v ní vystupuje neznámá $x$ a zároveň číslo $a$, jehož hodnotu neurčujeme. Řešením je diskuse o tom, kdy existuje jeden kořen, žádný nebo nekonečně mnoho.
\begin{definiceBox}
\textbf{Parametr v rovnici}. Pevně zvolené číslo, na němž závisí koeficienty nebo pravá strana rovnice. Při řešení hledáme neznámou $x$, nikoli hodnotu parametru.
\end{definiceBox}
\blok{Princip a odvození}
Nejprve uprav rovnici na $A(a)x=B(a)$. Než vydělíš výrazem $A(a)$, zvlášť prověř hodnoty parametru, pro které $A(a)=0$. V těchto případech se rovnice může změnit na pravdivou nebo nepravdivou konstantní rovnost.
\blok{Řešené příklady: od jednoduchého ke složitějšímu}
\begin{prikladReseniBox}[title={Příklad 1}]
Vyřešte $(a-2)x=a+1$ pro všechna $a\in\mathbb R$.
\tcblower
\begin{align*}
a&\ne2:\quad x=\frac{a+1}{a-2} \\
 a&=2:\quad 0\cdot x=3\;\text{-- spor}
\end{align*}

Pro $a\ne2$ jediný kořen $x=(a+1)/(a-2)$; pro $a=2$ žádný.
\end{prikladReseniBox}
\begin{prikladReseniBox}[title={Příklad 2}]
Vyřešte $(a-1)x=a-1$.
\tcblower
\begin{align*}
a&\ne1:\quad x=\frac{a-1}{a-1}=1 \\
 a&=1:\quad 0\cdot x=0
\end{align*}

Pro $a=1$ jsou řešením všechna reálná $x$.
\end{prikladReseniBox}
\poznamka{Jde o jedinou základní hodinu parametrických rovnic. Obtížnější víceparametrové úlohy nezařazuj jako povinné.}
\blok{Samostatné procvičování}
\begin{zadaniReseniBox}[title={Úloha 1}]
$(p+1)x=3$.
\tcblower
\begin{align*}
p&\ne-1:\quad x=\frac{3}{p+1} \\
 p&=-1:\quad 0\cdot x=3\;\text{-- spor}
\end{align*}
\end{zadaniReseniBox}
\begin{zadaniReseniBox}[title={Úloha 2}]
$(p-2)x=p-2$.
\tcblower
\begin{align*}
p&\ne2:\quad x=1 \\
 p&=2:\quad 0\cdot x=0\Rightarrow x\in\mathbb R
\end{align*}
\end{zadaniReseniBox}
\begin{zadaniReseniBox}[title={Úloha 3}]
$(p+3)x=0$.
\tcblower
\begin{align*}
p&\ne-3:\quad x=0 \\
 p&=-3:\quad 0\cdot x=0\Rightarrow x\in\mathbb R
\end{align*}
\end{zadaniReseniBox}
\blok{Volitelná rezerva}
\begin{zadaniReseniBox}[title={Úloha 1}]
Proč nelze z $ax=a$ automaticky napsat $x=1$?
\tcblower
\begin{align*}
a&\ne0:\quad x=1 \\
 a&=0:\quad 0=0\Rightarrow x\in\mathbb R
\end{align*}

Při $a=0$ má rovnice nekonečně mnoho řešení.
\end{zadaniReseniBox}
\blok{Domácí úkol}
\begin{zadaniReseniBox}[title={Domácí příprava}]
\begin{enumerate}[label=\alph*)]
\item
$(m-4)x=5$.
\item
$(m+2)x=m+2$.
\item
$mx=m+1$.
\item
$(a-3)x=2a-6$.
\item
$(a+1)x=a-2$.
\item
$(a-2)x=a+4$.
\end{enumerate}
\tcblower
\begin{enumerate}[label=\alph*)]
\item
\begin{align*}
m&\ne4:\quad x=\frac5{m-4} \\
 m&=4:\quad 0=5\;\text{-- spor}
\end{align*}
\item
\begin{align*}
m&\ne-2:\quad x=1 \\
 m&=-2:\quad 0=0\Rightarrow x\in\mathbb R
\end{align*}
\item
\begin{align*}
m&\ne0:\quad x=\frac{m+1}{m} \\
 m&=0:\quad 0=1\;\text{-- spor}
\end{align*}
\item
$\begin{aligned}
a\ne3&:\quad x=\frac{2(a-3)}{a-3}=2\\a=3&:\quad 0=0,\ \ K=\mathbb R
\end{aligned}$
\item
$\begin{aligned}
a\ne-1&:\quad x=\frac{a-2}{a+1}\\a=-1&:\quad 0=-3,\ \ K=\varnothing
\end{aligned}$
\item
$\begin{aligned}
a\ne2&:\quad x=\frac{a+4}{a-2}\\a=2&:\quad 0=6,\ \ K=\varnothing
\end{aligned}$
\end{enumerate}
\end{zadaniReseniBox}
\poznamka{Orientační časové rozvržení: 5 + 12 + 12 + 10 + 6 = 45 min. Přednost má pochopení jádra; rezerva není součástí časového součtu. Domácí úlohy jsou určeny na samostatnou práci, řešení zůstávají učiteli.}
\end{document}
